This study demonstrates that low-montage resting-state EEG can support patient-level prediction of remission outcomes in a home-based tDCS trial using an extremely sparse feature-selection framework. The strongest model achieved promising discrimination (ROC-AUC 0.816, MCC 0.571) while selecting only a small set of features, though the result does not survive correction for the full sweep of 150 configurations and should be considered exploratory.

Rather than relying on high-dimensional feature spaces, the results suggest that compact feature candidates centered on temporal entropy and frontal-temporal interactions may capture clinically relevant signal. These findings extend prior PSD-based analyses in the same cohort and support a convergent regional signal interpretation. However, the limited sample size ($n=21$), absence of an external validation cohort, and modest feature stability (Jaccard 0.190) indicate that these features should be treated as hypothesis-generating rather than as validated clinical biomarkers. Prospective validation in independent cohorts with pre-registered analysis plans is necessary before any clinical application can be considered.
